%% 303_Repraesentationale_Identitaet_De_ch.tex
%% FFGFT -- Dok. 303: Repräsentationale Identität über Rollenwechsel;
%%                    eine Governance-Ebene unter dem Claim-State-Ledger
%% Autor: Johann Pascher
%% Datum: Juli 2026
%% Sprache: Deutsch

\section*{Dok.~303 \quad Repräsentationale Identität über Rollenwechsel ---
eine Governance-Ebene unter dem Claim-State-Ledger}

Der Claim-State-Ledger fragt, ob ein Übergang \emph{verdient} ist: ob eine Aussage als
bewiesen, eingeschränkt oder blockiert gebucht werden darf. Diesem Test geht eine
frühere Frage voraus, die er nicht stellt --- ob das Objekt, das den Übergang durchläuft,
dabei seine \emph{repräsentationale Identität} bewahrt hat. Ein Rollenwechsel ist kein
Stärkewechsel. Dasselbe Symbol, dieselbe Zahl, dieselbe Struktur kann nacheinander
verschiedene repräsentationale Rollen besetzen --- mathematisches Objekt, Rechenergebnis,
Kandidat für eine physikalische Größe, Observable, ontologische Setzung ---, und zwischen
diesen Rollen verläuft die eigentliche, oft unbemerkte Verschiebung. Das vorliegende
Dokument benennt diese Ebene, gibt zwei Fallstudien (eine fremde, eine eigene), trennt die
Definition (die identifiziert) von der Governance (die bewahrt) und schlägt einen
Identitäts-Erhaltungs-Nachweis als dem Ledger vorgelagerte Ebene vor.

\section{Rollen, nicht Stärken}

Ein repräsentationales Objekt kann in einem Rahmen mindestens die folgenden Rollen
einnehmen:

\begin{itemize}
\item \emph{mathematisches Objekt} --- eine Größe innerhalb einer formalen Struktur, ohne
physikalische Zuschreibung;
\item \emph{Rechenkonstrukt / Ausgabe} --- das deterministische Resultat einer Prozedur;
\item \emph{Kandidat für eine physikalische Größe} --- ein Objekt, dem eine physikalische
Bedeutung \emph{vorgeschlagen}, aber noch nicht gesichert ist;
\item \emph{Observable} --- eine Größe, die an die Messung gekoppelt ist;
\item \emph{ontologische Setzung} --- eine Behauptung über das, was \emph{ist}, jenseits
von Ableitung und Falsifikation.
\end{itemize}

Diese Rollen sind nicht stärkere oder schwächere Versionen derselben Aussage; es sind
verschiedene \emph{Typen}. Der Claim-State-Ledger sortiert Aussagen nach ihrer
Belegstärke (bewiesen / eingeschränkt / blockiert); die Rollenfrage sortiert Objekte
nach ihrem repräsentationalen Typ. Beide Achsen sind nötig, und sie sind orthogonal.

\section{Der Fehlermodus: Rollendrift}

Der typische Fehler ist nicht die falsche Belegstärke, sondern die \emph{stille
Rollendrift}: dasselbe numerische oder mathematische Objekt wechselt seinen Typ, ohne
dass der Wechsel als solcher ausgewiesen wird. Eine Zahl, die als Gitterausgabe entsteht,
wird als \enquote{bare} physikalische Masse etikettiert, dann als Kandidat-Observable
geführt, schließlich als beobachtbarer Wert nach einer noch ausstehenden Korrektur
behandelt. Das sind vier repräsentationale Rollen desselben Zahlobjekts. Keine davon ist
an sich unzulässig; unzulässig ist nur, den Rollenwechsel als bloße Stärkung derselben
Aussage auszugeben. Genau an dieser Stelle tritt Uneinigkeit leise ein --- oft lange
bevor die Belegfrage überhaupt gestellt wird.

\section{Fallstudie I (fremd): eine Zahl in vier Rollen}

Ein numerisches Objekt aus einem diskreten Gittermodell durchlief in einer
IPI-Diskussion vier Rollen: zunächst \emph{Gitterausgabe} (ein deterministisches
Rechenergebnis), dann \emph{topologische Rohmasse} (eine physikalische Zuschreibung),
dann \emph{Kandidat} für ein beobachtbares Verhältnis, schließlich \emph{Observable nach
Abschirmung} (unter einem noch nicht konstruierten Reduktionsmechanismus). Der
Fortschritt der Diskussion bestand nicht darin, eine dieser Rollen zu widerlegen,
sondern darin, sie zu \emph{trennen}: sobald die Rollen auseinandergelegt waren, konnte
jede auf ihrem eigenen konstitutionellen Boden geprüft werden --- die Gitterausgabe als
bewiesen, die physikalische Zuordnung als Definition, das Observable als noch offener
Übergang. Die Trennung der Rollen war der eigentliche methodische Gewinn, unabhängig vom
Schicksal des Modells.

\section{Fallstudie II (eigen): die Informationseinheit}

Dieselbe Rollendrift trat bei der Ausarbeitung der vorliegenden Dokumentenserie an FFGFT
selbst auf. Die Frage nach der \enquote{minimalen Informationseinheit} wurde nacheinander
mit vier verschiedenen Objekten beantwortet: erst mit einer \emph{spektralen Eigenmode}
(einem Objekt des Massenoperators), dann mit einer \emph{Fläche} (der holographischen
Wahl eines Bit pro Flächenelement), dann mit einer \emph{Zählgröße} der Ordnung
$1/\xi$, schließlich mit dem \emph{topologischen Windungsquant} $\Delta w = 1$. Diese
vier sind nicht stärkere und schwächere Fassungen derselben Einheit; es sind vier
verschiedene repräsentationale Objekte, die nacheinander denselben Namen trugen. Der
Rollenwechsel blieb unbemerkt, bis die Identität durch eine \emph{vorgängige, explizite
Definition} fixiert wurde: Dok.~257 legt die Informationseinheit per Definition als
Windungsquant $\Delta w = 1$ fest ($\pi_1(T^4) \cong \mathbb{Z}^4$, skalenuniversell).
Erst dieser Identitätsanker machte die übrigen Rollen als das lesbar, was sie sind: die
Fläche als verlustbehaftete Projektion, das $\ln 2$ als thermodynamischer Schatten, die
Energie $\hbar c/L$ als skalenspezifische Projektion --- allesamt \emph{abgeleitete oder
projizierte} Erscheinungen der einen, fixierten Einheit, nicht konkurrierende
Definitionen.

\section{Definition identifiziert, Governance bewahrt}

Die Lehre aus beiden Fällen scheint zunächst zu sein, dass ein \emph{Identitätsanker}
die Identität bewahrt. Doch das vermengt zwei Fragen, die getrennt gehören --- und die
Vermengung ist genau die Rollendrift, die dieses Dokument beschreibt, nun auf das
Dokument selbst angewandt. Es sind zwei konstitutionelle Fragen:

\begin{itemize}
\item die \emph{deklarative} Frage --- \emph{Was ist dieses Objekt?} --- wird durch eine
\emph{Definition} beantwortet. Die Definition \emph{identifiziert}: sie setzt den
konstitutionellen Bezugspunkt, gegen den jeder weitere Anspruch überhaupt erst geprüft
werden kann. In FFGFT leistet das $\Delta w = 1$ (Dok.~257): sie legt fest, \emph{was}
die Informationseinheit ist.
\item die \emph{Governance}-Frage --- \emph{Unter welchen Operationen bleibt es dasselbe
Objekt?} --- ist davon verschieden. Eine Definition setzt die Anfangsreferenz; sie
\emph{bewahrt} die Identität nicht. Bewahrt wird sie erst durch die Governance über die
nachfolgenden Übergänge.
\end{itemize}

Kurz: \emph{die Definition identifiziert, die Governance bewahrt}. Der frühere Begriff
eines \enquote{Identitätsankers} trug beide Rollen unter einem Namen --- Definition
\emph{und} Erhaltungsmechanismus --- und ist damit selbst ein Fall der hier beschriebenen
Drift. Sauberer ist, ihn aufzulösen: die Definition ist der deklarierte Bezugspunkt; die
Governance ist eine eigene Ebene, die aufzeichnet, was mit diesem Bezug bei jedem
Übergang geschieht.

\section{Der Identitäts-Erhaltungs-Nachweis}

Eine bloß binäre Prüfung --- gleicher Typ vor und nach dem Übergang $\to$ Ledger;
verschiedener Typ $\to$ Substitution --- ist zu grob. Ein Übergang kann die deklarierte
Identität eines Objekts auf mehr als zwei Weisen behandeln, und die Governance-Ebene
führt darüber einen expliziten Nachweis: sie hält für jeden Übergang fest, ob er die
Identität

\begin{itemize}
\item \emph{bewahrt} --- dasselbe Objekt in derselben Rolle;
\item \emph{projiziert} --- eine skalen- oder koordinatenabhängige Erscheinung desselben
Objekts, verlustbehaftet und nicht fundamental;
\item \emph{transformiert} --- eine ableitbare Umformung, die die Referenz erhält;
\item \emph{substituiert} --- ein anderes Objekt unter demselben Namen (undeklariert ist
das die Rollendrift; deklariert braucht es eine eigene Lizenz);
\item \emph{überbrückt} --- eine gesicherte Abbildung auf ein Objekt eines anderen
Rahmens.
\end{itemize}

An FFGFTs eigener Informationseinheit gelesen: das Windungsquant $\Delta w = 1$ ist
\emph{bewahrt}; Energie ($\hbar c/L$) und Fläche sind \emph{Projektionen}; das $\ln 2$
ist eine \emph{Transformation} (der thermodynamische Schatten); die frühere Wanderung
Eigenmode $\to$ Fläche $\to 1/\xi$ war \emph{Substitution} unter beibehaltenem Namen, bis
die Definition (Dok.~257) die Referenz fixierte; und die Verifikationsrichtung
FFGFT$\to$HLV ist eine \emph{Brücke}. Der Nachweis entscheidet keine Physik; er stellt
nur sicher, dass bei jedem Schritt sichtbar bleibt, was mit der Identität geschieht ---
bevor der Ledger fragt, ob der Schritt verdient ist.

\section{Verhältnis zum Ledger}

Damit sind drei Ebenen getrennt. Die \emph{Definition} identifiziert das Objekt
(deklarativ). Die \emph{Governance} bewahrt seine Identität über Übergänge und hält im
Erhaltungs-Nachweis fest, ob ein Schritt bewahrt, projiziert, transformiert, substituiert
oder überbrückt. Der \emph{Claim-State-Ledger} bewertet, ob ein Übergang \emph{verdient}
ist (bewiesen / eingeschränkt / blockiert). Die Governance ist dem Ledger vorgelagert: sie
stellt sicher, dass über \emph{dasselbe} Objekt gesprochen wird, bevor über dessen
Belegstärke gesprochen wird. Alle drei Ebenen sind weitgehend unabhängig von der zugrunde
liegenden Physik --- sie betreffen die Buchführung der Repräsentation, nicht ihren Inhalt
--- und gewinnen an Bedeutung, sobald unabhängig entwickelte Rahmen aufeinandertreffen.
Ob diese Buchführung informell geführt wird oder durch einen systematischen Mechanismus,
ändert nichts an ihrer Wirkung: die Diskussion verschiebt sich vom Verteidigen von
Positionen zum Lokalisieren dessen, wo Annahmen sitzen, was abgeleitet ist und wo ein
Objekt seine Identität wechselt.

\section{Status}

Dieses Dokument ist ein methodischer Vorschlag, keine physikalische Ableitung. Die
FFGFT-spezifischen Aussagen --- die Informationseinheit als $\Delta w = 1$, die drei
Ebenen, die Projektionsnatur von Energie und Fläche --- sind korpusgestützt (Dok.~257,
301, 300); die Rollentaxonomie und die Identitätsprüfung sind konstitutionelle
Buchführung und als solche gebucht. Die zweite Fallstudie ist ausdrücklich
selbstkritisch: FFGFT bedurfte der hier vorgeschlagenen Disziplin selbst, und der
Identitätsanker $\Delta w = 1$ wurde erst nachträglich (Dok.~257) als solcher erkannt.
Der Anlass der Frage --- ein Einwand gegen Elementarzellen, die konstitutionelle Rahmung
des Übergangs-Problems und die repräsentationale Typisierung der beteiligten Objekte, wie
sie im IPI-Kreis (Takayuki Ueda, Stefaan Vossen, José Tomás Guevara Calderón) entwickelt
wurde --- wird dankend anerkannt. Die Unterscheidung zwischen deklarativer Identität
(Definition) und ihrer Erhaltung durch Governance sowie der fünfteilige
Erhaltungs-Nachweis (bewahren / projizieren / transformieren / substituieren /
überbrücken) wurden im Austausch mit Stefaan Vossen geschärft; sie sind gemeinsam
sichtbar gewordene konstitutionelle Sprache, kein Ersatz für und keine Übernahme eines
bestimmten Governance-Programms.
